\section{A uniform volume estimate}
\label{sec:volume}

Keep the background K\"ahler form $k$ fixed, and let $D$ be the
uniform bound from Lemma~\ref{ex:brody}:
\[
 d(f)=k\bigl(f'(0),f'(0)\bigr)\le D
\]
for every holomorphic disc in $X$.  We first extract a second consequence
of the extremal-disc argument.  Its constant will be independent of the
K\"ahler metric to which it is applied.

\begin{proposition}[Uniform lower bound for the volume form]
\label{prop:volume-lower-bound}
For every smooth K\"ahler form $h$ on $X$ satisfying
$R_h\ge-h$, one has
\begin{equation}
\label{eq:uniform-volume-lower-bound}
 h^n\ge c_{n,D}\,k^n,
 \qquad
 c_{n,D}=e^{-D}\left(\frac{n}{2n+D}\right)^n>0.
\end{equation}
\end{proposition}

\begin{proof}
Fix $x\in X$.  Use the maximizer from Proposition~\ref{ex:maximizer} with
$q=-h$ and $\delta=1$, so that its functional is
\[
 J(f)=d(f)-P_h(f)-A_h(f).
\]
The constant disc centered at $x$ has value zero.  Consequently a
maximizing disc $f$ satisfies
\begin{equation}
\label{eq:volume-extremal-energy}
 P_h(f)+A_h(f)\le d(f)\le D.
\end{equation}
Both terms on the left are nonnegative.  The assumed Ricci inequality
therefore gives
\[
 A_{R_h}(f)\ge-A_h(f)\ge-D,
 \qquad
 P_{R_h}(f)\ge-P_h(f)\ge-D.
\]

Let $B$ be the holomorphic frame along $f$ normalized to be
$h$-orthonormal on the boundary.  Since $\tr_h(-h)=-n$, the
maximum inequality, Equation~\eqref{eq:maximum-inequality}, yields
\[
 \tr\bigl(B(0)^*k(x)B(0)\bigr)
 \le 2n-A_{R_h}(f)\le2n+D.
\]
Meanwhile, the determinant identity,
Equation~\eqref{eq:determinant-identity}, gives
\[
 \det\bigl(B(0)^*h(x)B(0)\bigr)
 =\exp\bigl(P_{R_h}(f)\bigr)\ge e^{-D}.
\]
The arithmetic--geometric mean inequality for the eigenvalues of the
positive Hermitian matrix $B(0)^*k(x)B(0)$ gives
\[
 \det\bigl(B(0)^*k(x)B(0)\bigr)
 \le\left(\frac{2n+D}{n}\right)^n.
\]
Taking the quotient cancels the common factor $|\det B(0)|^2$:
\[
 \frac{h^n}{k^n}(x)
 =\frac{\det\bigl(B(0)^*h(x)B(0)\bigr)}
        {\det\bigl(B(0)^*k(x)B(0)\bigr)}
 \ge e^{-D}\left(\frac{n}{2n+D}\right)^n.
\]
The point $x$ was arbitrary, proving the assertion.
\end{proof}

The maximizing disc, its boundary regularity, and the size of its
holomorphic variation may all depend on $h$ and $x$.  The proof applies
the variational identities separately for each such choice.  None of
these auxiliary bounds enters the constant in
Equation~\eqref{eq:uniform-volume-lower-bound}.

\section{Positive top self-intersection}
\label{sec:positive-intersection}

Nefness allows us to approach the canonical class through K\"ahler
classes. We choose in each such class a metric to which the uniform
volume estimate applies. This Monge--Amp\`ere and cohomological-volume
strategy is also used by Wu and Yau
\cite[Proposition~8(i) and the proof of Lemma~6 in the authors' preprint]{WuYau2016}.
We use the following established form of the negative-sign complex
Monge--Amp\`ere theorem.  It is the Aubin--Yau existence theorem
\cite{Aubin1998,Yau1978}; the statement for an arbitrary smooth volume form is
also recorded explicitly in \cite[\S1, Equation~(1.1)]{Berman2019}.

\begin{theorem}[Negative-sign complex Monge--Amp\`ere theorem]
\label{thm:negative-cma}
Let $M$ be a compact K\"ahler manifold of dimension $n$, let $\gamma$
be a K\"ahler form on $M$, and let $\Omega$ be a smooth positive volume
form.  There is a smooth real function $w$ such that
\[
 \gamma+\ddc w>0,
 \qquad
 (\gamma+\ddc w)^n=e^w\Omega.
\]
\end{theorem}

There is no prescribed normalization of $\Omega$ in this statement.
The unknown includes its additive constant, which enters the
right-hand side through $e^w$.  In particular, the theorem is not
restricted to a K\"ahler form representing the canonical class.

\begin{proposition}
\label{prop:positive-top-intersection}
The canonical class has positive top self-intersection:
\[
 \int_X c_1(K_X)^n>0.
\]
\end{proposition}

\begin{proof}
Write
\[
 \alpha=[-R_k]=2\pi c_1(K_X).
\]
The nefness proved in Proposition~\ref{nef:nef} means that, for each
$t>0$, there is
a smooth real function $a_t$ for which
\[
 -R_k+\ddc a_t\ge-\frac t2 k.
\]
It follows that
\[
 \gamma_t=-R_k+tk+\ddc a_t\ge\frac t2 k>0
\]
is a K\"ahler form.  Apply Theorem~\ref{thm:negative-cma} with this
form and with the smooth positive volume form
$\Omega_t=e^{a_t}k^n$.  The resulting K\"ahler form
\[
 h_t=\gamma_t+\ddc w_t
\]
satisfies
\begin{equation}
\label{eq:twisted-cma-volume}
 h_t^n=e^{a_t+w_t}k^n.
\end{equation}
Taking $-\ddc$ of the logarithm of the ratio of the volume forms in
Equation~\eqref{eq:twisted-cma-volume} shows, with our Ricci convention, that
\begin{equation}
\label{eq:twisted-ricci-equation}
 R_{h_t}=R_k-\ddc(a_t+w_t)=-h_t+tk\ge-h_t.
\end{equation}
Proposition~\ref{prop:volume-lower-bound} therefore applies, with the
same constant for every $t>0$:
\[
 \int_Xh_t^n\ge c_{n,D}\int_Xk^n.
\]
On the other hand $[h_t]=\alpha+t[k]$, so closedness and Stokes'
theorem give
\[
 \int_Xh_t^n=\int_X(\alpha+t[k])^n.
\]
The expression on the right is a polynomial in $t$.  Passing to its
value at zero yields
\begin{equation}
\label{eq:positive-canonical-volume}
 (2\pi)^n\int_Xc_1(K_X)^n
 =\int_X\alpha^n
 \ge c_{n,D}\int_Xk^n>0.
\end{equation}
This passage uses only cohomology.  No convergence of the metrics
$h_t$, or uniform bounds on their potentials, is required.
\end{proof}

\section{Bigness, projectivity, and ampleness}
\label{sec:ampleness}

At this point $K_X$ is nef and has positive top self-intersection.
We first obtain enough sections to make $X$ Moishezon, then use its
K\"ahler structure to obtain projectivity. Only after that step do
we invoke the projective cone theorem to prove ampleness.
We give the positivity argument in the form needed here, including the
estimate that permits the use of holomorphic Morse inequalities.
For a real curvature form $\Theta$ of a smooth Hermitian line-bundle
metric, let $X(\Theta,\le1)$ denote the locus where $\Theta$ is
nondegenerate and has at most one negative eigenvalue.  The number of
negative eigenvalues is independent of the chosen background
Hermitian metric.

Demailly's strong holomorphic Morse inequalities
\cite[Theorem~0.1 and Equation~(0.7)]{Demailly1985} imply the following
criterion: on a compact complex manifold, if
\begin{equation}
\label{eq:morse-bigness-criterion}
 \int_{X(\Theta,\le1)}\Theta^n>0,
\end{equation}
then the line bundle is big and the manifold is Moishezon
\cite[Theorem~0.8(a)]{Demailly1985}.  Here ``big'' means that the
line bundle has Kodaira dimension $n$; equivalently its spaces of
sections have growth of order $m^n$ along sufficiently divisible
tensor powers.  A compact complex manifold is Moishezon when its field
of meromorphic functions has transcendence degree $n$.

\begin{lemma}[Nefness and positive intersection imply bigness]
\label{lem:nef-big}
Let $L$ be a nef holomorphic line bundle on a compact connected K\"ahler
manifold $(X,k)$ of dimension $n\ge1$.  If
$\int_Xc_1(L)^n>0$, then $L$ is big and $X$ is Moishezon.
\end{lemma}

\begin{proof}
Put $\beta=2\pi c_1(L)$ and $V=\int_X\beta^n>0$.
For $0<\eps\le1$, choose a smooth metric on $L$ with real curvature
$\Theta_\eps\ge-\eps k$.  Thus
\[
 P_\eps=\Theta_\eps+2\eps k
\]
is positive definite.  At a point where $\Theta_\eps$ has a negative
eigenvalue, write its eigenvalues relative to $k$ as
$\lambda_1,\ldots,\lambda_n$ and choose $j$ with $\lambda_j<0$.
Then
\[
 |\lambda_j|\le\eps,
 \qquad
 |\lambda_i|\le\lambda_i+2\eps\quad(i\ne j).
\]
Consequently
\[
 \left|\prod_{i=1}^n\lambda_i\right|
 \le\eps\prod_{i\ne j}(\lambda_i+2\eps)
 \le\eps\sum_{j=1}^n\prod_{i\ne j}(\lambda_i+2\eps).
\]
In terms of volume forms this is
\begin{equation}
\label{eq:negative-index-mass}
 |\Theta_\eps^n|\le n\eps\,P_\eps^{n-1}\wedge k
\end{equation}
on the entire locus with a negative eigenvalue.  The right-hand side
has integral
\[
 n\eps\int_X(\beta+2\eps[k])^{n-1}\wedge[k]=O(\eps).
\]
The implied constant is independent of the chosen metric: the
remaining cohomological integral is a fixed polynomial in $\eps$,
bounded for $0\le\eps\le1$.

The degenerate locus contributes zero to $\Theta_\eps^n$.  Removing
the loci of index at least two from its total integral therefore
changes that integral by at most $O(\eps)$ in absolute value.  Since
$\int_X\Theta_\eps^n=V$, we obtain
\[
 \int_{X(\Theta_\eps,\le1)}\Theta_\eps^n
 \ge V-O(\eps)>0
\]
for a sufficiently small fixed $\eps$.  The criterion in
Equation~\eqref{eq:morse-bigness-criterion} applies to this one smooth metric
and proves the assertion.  In dimension one the locus of index at
least two is empty; the same argument, with $P_\eps^0=1$, applies
without modification.
\end{proof}

Apply Lemma~\ref{lem:nef-big} to $L=K_X$, using Proposition~\ref{nef:nef} and
Proposition~\ref{prop:positive-top-intersection}.  We conclude that
$K_X$ is big and $X$ is Moishezon.  The following established
projectivity theorem now applies.

\begin{theorem}[Moishezon's projectivity theorem]
\label{thm:moishezon-projectivity}
A compact Moishezon manifold that admits a K\"ahler metric is
projective.
\end{theorem}

For a primary source proving a stronger form of this theorem, see
Theorem~6 in the author's preprint of \cite{Namikawa2002}: it permits
Moishezon spaces with
$1$-rational singularities.  A smooth manifold satisfies that
hypothesis.  Thus $X$ is now a smooth projective variety, and we may
apply the projective cone theorem.

The following standard implication is recorded in
\cite[Lemma~2.1]{DiverioTrapani2019}. We give the argument through an
effective perturbation of the canonical divisor and the log cone
theorem.

\begin{proposition}[Bigness in the absence of rational curves]
\label{prop:no-rational-ample}
Let $X$ be a smooth projective variety over $\C$.  If $K_X$ is big
and $X$ contains no rational curves, then $K_X$ is ample.
\end{proposition}

\begin{proof}
Kodaira's lemma gives an integer $m>0$, an ample Cartier divisor $H$,
and an effective divisor $E$ such that
\begin{equation}
\label{eq:kodaira-decomposition}
 mK_X\sim H+E.
\end{equation}
For the needed form of this lemma, choose a smooth very ample
hyperplane section $H$. Bigness gives $h^0(X,mK_X)\ge cm^n$ along
sufficiently divisible powers, whereas the weak holomorphic Morse
inequality gives $h^0(H,mK_X|_H)=O(m^{n-1})$
\cite[Theorem~0.1(a)]{Demailly1985}. The restriction sequence therefore
gives a nonzero section of $mK_X-H$ for some $m$, producing $E$.
For $n=1$, $H$ is a finite set and the same estimate is $O(1)$.

Choose a sufficiently small positive rational number $\eta$ so that
$(X,\eta E)$ is Kawamata log terminal.  This choice is possible even
when $E$ is singular or nonreduced. Choose a log resolution
$\pi:Y\to X$ of $(X,E)$, whose existence follows from
\cite[Theorem~35]{Kollar2007}, and write
\[
 K_Y-\pi^*K_X=\sum_i a_iF_i,
 \qquad \pi^*E=\widetilde E+\sum_i b_iF_i,
\]
where $F_i$ are exceptional prime divisors and $\widetilde E$ is the
strict transform with its multiplicities. The discrepancy criterion
for a Kawamata log terminal pair is given in
\cite[\S4.4]{Fujino2011}. The exceptional discrepancies have the form
$a_i-\eta b_i$, with $a_i\ge0$ because $X$ is smooth and $b_i\ge0$
because $E$ is effective.  There are finitely many such coefficients.
Taking $\eta>0$ sufficiently small keeps each discrepancy greater
than $-1$ and each strict-transform boundary coefficient less than
one, which is the required condition.

The log cone theorem says that a projective Kawamata log terminal
pair $(X,\Delta)$ over $\C$ with effective rational boundary has a
rational curve of negative $(K_X+\Delta)$-degree whenever
$K_X+\Delta$ is not nef.  This follows, for example, from the cone
decomposition and rational-curve assertion in
\cite[Theorem~1.1, especially part~(5)]{Fujino2011}; for a Kawamata
log terminal pair the non-log-canonical locus in that statement is
empty.  The theorem has no restriction on the dimension.

Since $X$ has no rational curves, it follows that
\[
 N=K_X+\eta E
\]
is nef, in the numerical sense that its degree on every curve is
nonnegative. Combining this with Equation~\eqref{eq:kodaira-decomposition} gives
\[
 (1+\eta m)K_X\sim_{\mathbb Q}N+\eta H.
\]
The sum of a nef rational divisor and an ample rational divisor on a
projective variety is ample by Kleiman's ampleness criterion
\cite[Definition~4.9 and Theorem~4.10]{Fujino2011}: it is positive
on every nonzero element of the closed cone of curves. Hence
$(1+\eta m)K_X$ is ample. Clearing the
positive rational coefficient shows that a positive integral multiple
of $K_X$ is ample, and therefore $K_X$ itself is ample.
\end{proof}

\begin{proof}[Completion of the proof of Theorem~\ref{thm:main}]
A rational curve in $X$ would give a nonconstant holomorphic map
$\mathbb P^1\to X$ through its normalization.  Its restriction to
$\C\subset\mathbb P^1$ would still be nonconstant, by the identity
theorem.  The hypothesis therefore excludes rational curves.
We have proved that $X$ is projective and $K_X$ is big, so
Proposition~\ref{prop:no-rational-ample} proves that $K_X$ is ample.
Equivalently, a sufficiently high positive tensor power of $K_X$ is
very ample and its global holomorphic sections define a holomorphic
embedding into projective space.  This is the assertion in every
positive complex dimension.
\end{proof}
